%! Probability Models — Unit 3, Lesson 3.2 — Warm-Up
% ONE PAGE, blank and key. Three items at 12pt. Do not add a fourth item.
% GRADUAL RELEASE: fractions, decimals, percents and simplifying, on the numbers the notes reopen:
%   Item 1: 28/40 = 7/10 (problem 9, P(taco or lemonade)); 12/40 = 3/10 (problem 7, neither);
%           10/40 = 1/4 (problem 8, only lemonade); 80/120 = 2/3 (problem 6, P(lemonade)).
%   Item 2: 3/8 = 0.375 = 37.5% (problem 1, the free-drink sectors); 5/8 = 0.625 = 62.5%
%           (problem 3, the complement); 7/10 = 0.7 = 70% (problem 9).
%   Item 3: (3/8)(200) = 75 (problem 2, the expected count); 1 - 3/8 = 5/8 (problem 3).
\documentclass[12pt]{article}
\usepackage{saar-article}
\usepackage{saar-boxes}
\newcommand{\TallMath}[1]{$\displaystyle #1\rule[-1.4em]{0pt}{3.2em}$}
% Word bank strip for every fill-in cluster. Defined per document; the key inherits it and
% prints the same strip, so heights match.
\newcommand{\wordbank}[1]{\par\vspace{2pt}\noindent\fcolorbox{goldacc}{goldbg}{\parbox{\dimexpr\linewidth-2\fboxsep-2\fboxrule\relax}{\small\textbf{Word bank:} #1}}\par\vspace{4pt}}
% Handwriting room inside every \begin{work} block. Moves the blank and the
% key together, so the two can never drift apart.
\setlength{\workrowsep}{24pt}

\begin{document}

\pageheader{Unit 3, Lesson 3.2}{Warm-Up}
% NAMESTRIP: no \namedateperiod here. The student writes their name once, on the
% cover this component is stapled behind.

\vspace{0.15in}

\boxguard
\begin{notesbox}{Spiral Review --- Fractions, Decimals, and Percents}
\normalsize
A probability is written as a fraction, a decimal, or a percent. Today's notes use these exact
numbers.

\begin{enumerate}[leftmargin=1.8em, itemsep=20pt, topsep=10pt]

  \item Simplify each fraction.

        \vspace{4pt}
        \renewcommand{\arraystretch}{2.6}
        \begin{tabularx}{\linewidth}{@{} Y Y Y Y @{}}
        \toprule
        $\dfrac{28}{40}$ & $\dfrac{12}{40}$ & $\dfrac{10}{40}$ & $\dfrac{80}{120}$ \\
        \midrule
        \blank{2.2cm} & \blank{2.2cm} & \blank{2.2cm} & \blank{2.2cm} \\
        \bottomrule
        \end{tabularx}

  \item Write each fraction as a decimal and as a percent.
        \wordbank{0.375 \;\textbullet\; 37.5\% \;\textbullet\; 0.625 \;\textbullet\; 62.5\% \;\textbullet\; 0.7 \;\textbullet\; 70\%}

        \vspace{2pt}
        \renewcommand{\arraystretch}{2.2}
        \begin{tabularx}{\linewidth}{@{} X p{3.4cm} p{3.4cm} @{}}
        \toprule
        \textbf{Fraction} & \textbf{Decimal} & \textbf{Percent} \\
        \midrule
        $3/8$  & \blank{3cm} & \blank{3cm} \\
        $5/8$  & \blank{3cm} & \blank{3cm} \\
        $7/10$ & \blank{3cm} & \blank{3cm} \\
        \bottomrule
        \end{tabularx}

  \item Find $\dfrac{3}{8}$ of $200$. Then find what is left when you take $\dfrac{3}{8}$ away
        from one whole.

        \begin{work}
          \tfrac{3}{8} \cdot 200 &= 600 \div 8 = 75 \\
          1 - \tfrac{3}{8} &= \tfrac{8}{8} - \tfrac{3}{8} = \tfrac{5}{8}
        \end{work}

        \vspace{6pt}
        $\dfrac{3}{8}$ of $200$: \blank{2.4cm} \qquad\qquad What is left: \blank{2.4cm}

\end{enumerate}
\end{notesbox}

\end{document}
